\usepackage{float} % Improved interface for floating objects
\usepackage[greek,english]{babel}
\usepackage{ragged2e} % For text justification
% minitoc breaks the IEEEtran cls
% \usepackage{minitoc} % Allows you to include contents page in chapters if you want
\usepackage[table]{xcolor} % Advanced color handling with table option for coloring tables
\usepackage[backend=biber,
            style=ieee,
            sorting=nyt, %maxbibnames=99,
            doi=true,
            url=true,
            giveninits=true]{biblatex}
\usepackage{amsmath,amssymb,amsfonts} % AMS mathematical facilities and symbols
\usepackage{bm} % Access bold math symbols
\usepackage{graphicx} % Enhanced support for graphics
\usepackage{textcomp} % Additional text symbols
\usepackage{hyperref} % Extensive support for hypertext in LaTeX
\usepackage{tikz} % Creation of diagrams and pictures using PGF
\usetikzlibrary{positioning, shapes.multipart, arrows.meta}
\usepackage{mathdots} % Commands to produce dots in math that respect font size
\usepackage{yhmath} % Extended maths fonts for LaTeX
\usepackage{cancel} % Place lines through maths formulae
\usepackage{siunitx} % A comprehensive (SI) units package
\usepackage{array} % Extending the array and tabular environments
\usepackage{multirow} % Create tabular cells spanning multiple rows
\usepackage{gensymb} % Generic symbols for both text and math mode
\usepackage{tabularx} % Tabulars with adjustable-width columns
\usepackage{xurl} % Allow URL breaks at any alphanumerical character
\usepackage{listings} % Typeset source code listings using LaTeX
\usepackage{minted} % Highlighted source code for LaTeX
\usepackage{pgf} % Create PostScript and PDF graphics
\usepackage{pgf-pie} % Draw pie charts, part of the pgf bundle
\usepackage{tikz-cd} % Create commutative diagrams with TikZ
\usepackage[en-US]{datetime2} % Formats for dates, times and time zones
\usepackage{pgfplotstable} % Loads, rounds, formats and postprocesses numerical tables
\usepackage{booktabs} % Publication quality tables in LaTeX
\usepackage{pgfplots} % Create normal/logarithmic plots in two and three dimensions
\usepackage{arydshln} % Dash-Lined hline and vline in array and tabular
\usepackage[most]{tcolorbox} % Colored boxes, for LaTeX examples and theorems, etc.
\usepackage{pgf-umlsd} % Draw UML Sequence Diagrams
\usepackage{enumitem} % Modify the spacing between items
\usepackage{graphicx}
\usepackage{subcaption}
\usepackage{fbox}
\usepackage{setspace}
\usepackage{pifont} % For \ding{} symbols
\usepackage{ifthen}
%\usepackage{worldflags}
\usepackage{pst-flags}
\usepackage[expansion=true, kerning=true]{microtype} % widows

%\newsavebox{\USflagbox}
%\sbox{\USflagbox}{\worldflag[width=2mm]{US}}

% Define the 'definition' environment
%\newtheorem{definition}{Definition}[section]

\newtcbtheorem[number within=section]{theorem}{Theorem}{%
    enhanced,
    colback=gray!5,  % Light grey background
    colframe=gray!40, % Darker grey frame
    arc=4mm,         % Rounded corners
    boxrule=1pt,     % Frame thickness
    fonttitle=\bfseries,
    separator sign={.}
}{th}

%-----------------------------------------------------------------------------------------------
% Theorem style
\usepackage[ntheorem]{mdframed}
\mdfsetup{%
topline=false,
rightline=false,
bottomline=false,
linewidth=3pt,
innerleftmargin=15pt,
innerrightmargin=0pt,
skipabove=\baselineskip,
skipabove=1.2\baselineskip,
}

\newtheorem{mdbeispiel}{Beispiel}[section]
\newtheorem{mdspiele}{Theorem}[section]
\newenvironment{beispiel}[1][]%
   {\begin{mdframed}[linecolor=blue]\begin{mdbeispiel}[#1]}
   {\end{mdbeispiel}\end{mdframed}}
\newenvironment{Theorem}[1][]%
   {\begin{mdframed}[linecolor=yellow]\begin{mdspiele}[#1]}
   {\end{mdspiele}\end{mdframed}}

\newtheorem{conj}{Conjecture}[section]
\newenvironment{conjecture}[1][]%
   {\begin{mdframed}[linecolor=red]\begin{conj}[#1]}
   {\end{conj}\end{mdframed}}
   
\newtheorem{mddef}{Definition}[section]
\newenvironment{definition}[1][]%
   {\begin{mdframed}[linecolor=lightgray]\begin{mddef}[#1]}
   {\end{mddef}\end{mdframed}}

   
%-----------------------------------------------------------------------------------------------

% Define theorem-like environments
% \theoremstyle{plain}
% \newtheorem{theorem}{Theorem}[section]
% \newtheorem{lemma}[theorem]{Lemma}
% \newtheorem{proposition}[theorem]{Proposition}
% \newtheorem{corollary}[theorem]{Corollary}

% \theoremstyle{definition}
% \newtheorem{definition}[theorem]{Definition}
% \newtheorem{axiom}[theorem]{Axiom}
% \newtheorem{example}[theorem]{Example}
% 
% \theoremstyle{remark}
% \newtheorem{remark}[theorem]{Remark}
% \newtheorem{note}[theorem]{Note}

%-----------------------------------------------------------------------------------------------
%   Algorithms

%%  Confusing state of the algorithm packages:
%%  algorithm - float wrapper for algorithms.
%%  algorithmic - first algorithm typesetting environment.
%%  algorithmicx - second algorithm typesetting environment.
%%  algpseudocode - layout for algorithmicx.
%%  algorithm2e - third algorithm typesetting environment.

\usepackage[linesnumbered,ruled,vlined]{algorithm2e}
\DontPrintSemicolon
%-----------------------------------------------------------------------------------------------

%   Font-sizes

%   \tiny
%   \scriptsize
%   \footnotesize
%   \small
%   \normalsize	F-normalsize.png
%   \large	F-large.png
%   \Large	F-large2.png
%   \LARGE	F-large3.png
%   \huge	F-huge.png
%   \Huge	F-huge2.png

%-----------------------------------------------------------------------------------------------
%   Define pseudocode formatting
\renewcommand{\KwSty}[1]{\textnormal{\textcolor{blue!90!black}{\ttfamily\bfseries #1}}\unskip}
\renewcommand{\ArgSty}[1]{\textnormal{\ttfamily #1}\unskip}
\SetKwComment{Comment}{\color{green!50!black}// }{}
\renewcommand{\CommentSty}[1]{\textnormal{\ttfamily\color{green!50!black}#1}\unskip}
\newcommand{\assign}{\leftarrow}
\newcommand{\var}{\texttt}
\newcommand{\FuncCall}[2]{\texttt{\bfseries #1(#2)}}
\SetKwProg{Function}{function}{}{}
\renewcommand{\ProgSty}[1]{\texttt{\bfseries #1}}

%\def\BibTeX{{\rm B\kern-.05em{\sc i\kern-.025em b}\kern-.08em
%    T\kern-.1667em\lower.7ex\hbox{E}\kern-.125emX}}

\lstset{basicstyle=\ttfamily, keywordstyle=\bfseries}
%-----------------------------------------------------------------------------------------------
\pgfplotsset{compat=1.18}
%-----------------------------------------------------------------------------------------------
%   Colors
\definecolor{mygreen}{RGB}{28,172,0} % color values Red, Green, Blue
\definecolor{mylilas}{RGB}{170,55,241}
\definecolor{lightgreen}{HTML}{A1EEBD}
\definecolor{darkgreen}{HTML}{3f6e50}
\definecolor{pieone}{HTML}{FAE5D3}%light color
\definecolor{pietwo}{HTML}{A04000}
\definecolor{xLightGreen}{HTML}{A1EEBD}
\definecolor{xLightYellow}{HTML}{F6F7C4}
\definecolor{xLightRed}{HTML}{F6D6D6}
\definecolor{xLightBlue}{HTML}{7BD3EA}
\definecolor{deepblue}{rgb}{0,0,0.5}
\definecolor{deepred}{rgb}{0.6,0,0}
\definecolor{deepgreen}{rgb}{0,0.5,0}

% h Place the float here, i.e., approximately at the same point it occurs in the source text (however, not exactly at the spot)
% t Position at the top of the page.
% b Position at the bottom of the page.
% p Put on a special page for floats only.
% ! Override internal parameters LaTeX uses for determining "good" float positions.
% H Places the float at precisely the location in the LATEX code. Requires the float package. This is somewhat equivalent to h!.
% 

\newcommand{\DrawPercentageBar}[1]{%
  \begin{tikzpicture}
    \fill[color=black]   (0.0 , 0.0) rectangle (#1*7ex , 1.75ex );
    \fill[color=gray] (#1*7ex  , 0.0) rectangle (7.0ex, 1.75ex);
  \end{tikzpicture}%
}

\newcommand{\DrawLengthBar}[1]{%
  \begin{tikzpicture}
    \fill[color=lightgreen]   (0.0 , 0.0) rectangle (#1*0.8ex , 1.75ex );
  \end{tikzpicture}%
}

\newcommand{\DrawLengthBarThick}[1]{%
  \begin{tikzpicture}
    \fill[color=lightgreen]   (0.0 , 0.0) rectangle (#1*0.8ex , 2ex );
  \end{tikzpicture}%
}
\renewcommand\tabularxcolumn[1]{m{#1}}% for vertical centering text in X column

\newcommand{\DrawSmallPieChart}[1]{%
    \begin{tikzpicture}[draw opacity=0]
        \node at (3.8ex,-0.0ex) { \textcolor{darkgreen}{\footnotesize #1 } };
        \pie[scale=0.055, color={darkgreen, lightgreen}, before number=\phantom, after number=,
        style =very thin,
        /tikz/nodes={opacity=0,overlay}% <-
]{#1/, 100-#1/}
    \end{tikzpicture}
}

\newcommand{\Password}[1]{%
  \begin{tikzpicture}
    \small
    \node[draw, rectangle, rounded corners=4pt, inner sep=3pt] {\texttt{#1}};
  \end{tikzpicture}%
}

\usetikzlibrary{fadings}
\usetikzlibrary{patterns}
\usetikzlibrary{shadows.blur}
\def\BibTeX{{\rm B\kern-.05em{\sc i\kern-.025em b}\kern-.08em
    T\kern-.1667em\lower.7ex\hbox{E}\kern-.125emX}}

% Requires arydshln package
\setlength\dashlinedash{0.2pt}
\setlength\dashlinegap{2pt}
\setlength\arrayrulewidth{0.1pt}
%-----------------------------------------------------------------------------------------------
%% Maths stuff
\DeclareMathOperator*{\di}{\mathrm{d}\!}
\def\at{
  \left.
  \vphantom{\int}
  \right|
}



%-----------------------------------------------------------------------------------------------
%% Grey note box with red heading
\newcommand{\noteBox}[3][16cm]{%
    \begin{center}
    \begin{tikzpicture}
    \def\gridwidth{#1}% Optional width parameter (default 16cm)
    \node[draw=gray!30, fill=gray!3, text width=\gridwidth-20pt, 
          inner sep=10pt, rounded corners=2pt] (box) {
        \begin{minipage}{\linewidth}
        \centering
        \normalsize\textbf{\textcolor{purple}{#2}}% Title parameter
        
        \vspace{0.5em}
        \normalsize\setstretch{1.2}\justifying
        #3% Content parameter
        \end{minipage}
    };
    \end{tikzpicture}
    \end{center}
}



\newcommand{\noteBoxNormal}[3][16cm]{%
    \begin{center}
    \begin{tikzpicture}
    \def\gridwidth{#1}% Optional width parameter (default 16cm)
    \node[draw=gray!30, fill=gray!3, text width=\gridwidth-20pt, 
          inner sep=10pt, rounded corners=2pt] (box) {
        \begin{minipage}{\linewidth}
        \centering
        \textbf{\textcolor{purple}{#2}}% Title parameter
        
        \vspace{0.5em}
        \normalsize\justifying
        #3% Content parameter
        \end{minipage}
    };
    \end{tikzpicture}
    \end{center}
}

\newcommand{\noteBoxGrey}[3][\textwidth-20pt]{%
    \begin{center}
    \begin{tikzpicture}
    \def\gridwidth{#1}% Optional width parameter (default 16cm)
    \node[fill=gray!10, text width=\gridwidth-20pt, 
          inner sep=20pt, rounded corners=4pt] (box) {
        \begin{minipage}{\linewidth}
        \centering
        \large\textbf{\textcolor{purple}{#2}}% Title parameter

        \vspace{0.5em}
        \small\setstretch{1}\justifying
        #3% Content parameter
        \end{minipage}
    };
    \end{tikzpicture}
    \end{center}
}

\newcommand{\bibBox}[3][\textwidth-20pt]{%
    \begin{center}
    \begin{tikzpicture}
    \def\gridwidth{#1}% Optional width parameter (default 16cm)
    \node[fill=gray!10, text width=\gridwidth-20pt, 
          inner sep=5pt, rounded corners=4pt] (box) {
        \begin{minipage}{\linewidth}
        \centering
        \small\textbf{\textcolor{purple}{#2}}% Title parameter
        \hspace{0.2cm}
        \small\setstretch{1}\justifying
        #3% Content parameter
        \end{minipage}
    };
    \end{tikzpicture}
    \end{center}
}



%-----------------------------------------------------------------------------------------------
% Fancy text, e.g. password in a rounded rectangle and a diff colour
\newcommand{\password}[2][violet!30]{%
  \tikz[baseline=(textnode.base)]{
    \node[
      draw=#1,             % Outline color (default: black)
      text=violet,
      rounded corners=3.5pt, % Corner radius
      thin,                % Outline thickness
      inner sep=2pt,       % Padding between text and rectangle
      anchor=base          % Aligns with baseline
    ] (textnode) {\texttt{#2}};
  }%
}


\newcommand{\guessmodel}[2][deepblue!10]{%
  \tikz[baseline=(textnode.base)]{
    \node[
      draw=#1,             % Outline color (default: black)
      text=deepblue,
      rounded corners=3.5pt, % Corner radius
      thin,                % Outline thickness
      inner sep=2pt,       % Padding between text and rectangle
      anchor=base          % Aligns with baseline
    ] (textnode) {\texttt{#2}};
  }%
}

% Create a san-serif style in math mode
\newcommand{\fun}[1]{\ensuremath{\mathsf{#1}}}
\newcommand{\srcOwnFig}{Source: Own figure.\xspace}  % Shortcut for own figure.
\newcommand{\etal}{\textit{et al.}\xspace}
\newcommand{\passwordNinjaURL}{\url{https://github.com/sirvan3tr/passwordninja}\xspace}
%-----------------------------------------------------------------------------------------------
%% Thesis specific
\newcommand{\hmail}{\fun{hotmail}\xspace}
\newcommand{\phpbb}{\fun{phpbb}\xspace}
\newcommand{\lizardsquad}{\fun{LizardSquad}\,\cite{miessler_lizardsquad_2015}\xspace}

\newcommand{\lizardsquadnocite}{\fun{LizardSquad}\xspace}
\newcommand{\pwdninja}{\fun{passwordNinja}\xspace}

\newcommand{\hashcat}{Hashcat~\cite{steube_hashcat_2023}\xspace}
\newcommand{\jtr}{JtR~\cite{JtR}\xspace}

\newcommand{\markovngram}{\guessmodel{Markov $n$-gram}~\cite{narayanan_fast_2005}\xspace}
\newcommand{\pcfg}{\guessmodel{PCFG}~\cite{weir_password_2009}\xspace}
\newcommand{\zxcvbn}{\guessmodel{zxcvbn}~\cite{zxcvbn2016}\xspace}
\newcommand{\passgan}{\guessmodel{PassGAN}~\cite{hitaj_passgan_2019}\xspace}
\newcommand{\cwaecpg}{\guessmodel{CWAE-CPG}~\cite{pasquini_improving_2021}\xspace}
\newcommand{\cwaegan}{\guessmodel{CWAE-GAN}~\cite{pasquini_improving_2021}\xspace}
\newcommand{\omen}{\guessmodel{OMEN}~\cite{durmuth_omen_2015}\xspace}
\newcommand{\omecdn}{\guessmodel{OMECDN}~\cite{jiang_omecdn_2022}\xspace}
\newcommand{\rfguess}{\guessmodel{RFGuess}~\cite{wang_password_2023}\xspace}
\newcommand{\gnpassgan}{\guessmodel{GNPassGAN}~\cite{yu_gnpassgan_2022}\xspace}
\newcommand{\adams}{\guessmodel{AdaMs}~\cite{pasquini_reducing_2021}\xspace}
\newcommand{\fla}{\guessmodel{FLA}~\cite{melicher_fast_2016}\xspace}
\newcommand{\gpass}{\guessmodel{G-Pass}~\cite{zhou_password_2022}\xspace}
\newcommand{\redpack}{\guessmodel{REDPACK}~\cite{redpackgan2020}\xspace}
\newcommand{\wordpcfg}{\guessmodel{WordPCFG}~\cite{wordpcfg_2021}\xspace}
\newcommand{\passtcn}{\guessmodel{PassTCN-PPLL}~\cite{passtcn_ppll_2022}\xspace}
\newcommand{\personalpcfg}{\guessmodel{Personal-PCFG}~\cite{li_study_2016}\xspace}
\newcommand{\passbert}{\guessmodel{PassBERT}~\cite{xu_improving_2023}\xspace}

\newcommand{\markovngramnocite}{\guessmodel{Markov $n$-gram}}
\newcommand{\pcfgnocite}{\guessmodel{PCFG}}
\newcommand{\zxcvbnnocite}{\guessmodel{zxcvbn}}
\newcommand{\passgannocite}{\guessmodel{PassGAN}}
\newcommand{\cwaecpgnocite}{\guessmodel{CWAE-CPG}}
\newcommand{\cwaegannocite}{\guessmodel{CWAE-GAN}}
\newcommand{\omennocite}{\guessmodel{OMEN}}
\newcommand{\omecdnnocite}{\guessmodel{OMECDN}}
\newcommand{\rfguessnocite}{\guessmodel{RFGuess}}
\newcommand{\gnpassgannocite}{\guessmodel{GNPassGAN}}
\newcommand{\adamsnocite}{\guessmodel{AdaMs}}
\newcommand{\flanocite}{\guessmodel{FLA}}
\newcommand{\gpassnocite}{\guessmodel{G-Pass}}
\newcommand{\redpacknocite}{\guessmodel{REDPACK}}
\newcommand{\wordpcfgnocite}{\guessmodel{WordPCFG}}
\newcommand{\passtcnnocite}{\guessmodel{PassTCN-PPLL}}
\newcommand{\personalpcfgnocite}{\guessmodel{Personal-PCFG}}
\newcommand{\passbertnocite}{\guessmodel{PassBERT}\xspace}

\newcommand{\deeid}{\fun{deeID}\xspace}
\newcommand{\formless}{\fun{FormL3SS}\xspace}


%-----------------------------------------------------------------------------------------------
%% Key derivation functions
\newcommand{\bcrypt}{\texttt{Bcrypt}~\cite{provos_future-adaptable_1999}}
\newcommand{\pbkdf}{\texttt{PBKDF2}~\cite{kaliski_pkcs_2000}}
\newcommand{\spkdf}{\texttt{SP800-108}~\cite{chen_recommendation_2009}}
\newcommand{\scrypt}{\texttt{Scrypt}~\cite{percival_stronger_2009}}
\newcommand{\hkdf}{\texttt{HKDF}~\cite{krawczyk_hmac-based_2010}}
\newcommand{\sakke}{\texttt{SAKKE}~\cite{groves_sakke_2012}}
\newcommand{\catena}{\texttt{Catena}~\cite{forler_catena_2013}}
\newcommand{\sphinx}{\texttt{SPHINX}~\cite{sphinx_2015}}
\newcommand{\makwa}{\texttt{Makwa}~\cite{pornin_makwa_2015}}
\newcommand{\yescrypt}{\texttt{yescrypt}~\cite{neves_yescrypt_2015}}
\newcommand{\argontwo}{\texttt{Argon2}~\cite{biryukov_argon2_2016}}
\newcommand{\balloon}{\texttt{Balloon}~\cite{boneh_balloon_2016}}
\newcommand{\lyratwo}{\texttt{Lyra2}~\cite{andrade_lyra2_2016}}
\newcommand{\mixkey}{\texttt{MixKey}~\cite{perrin_mixkey_2017}}

%% Company names
\newcommand{\ha}{Hanna\;Andersson\xspace}

% Define colors
\definecolor{lightgray}{RGB}{245,245,245}
\definecolor{midgray}{RGB}{225,225,225}
\definecolor{recommend}{RGB}{144,238,144}
\definecolor{discourage}{RGB}{255,182,193}

% Custom symbols for recommendations
\newcommand{\recommended}{\textcolor{green}{\faCheckCircle}}
\newcommand{\discouraged}{\textcolor{red}{\faTimesCircle}}
\newcommand{\optional}{\textcolor{orange}{\faQuestionCircle}}
\newcommand{\required}{\textcolor{blue}{\faExclamationCircle}}
\newcommand{\secure}{\textcolor{green}{\faLock}}
\newcommand{\manager}{\faKey}
\newcommand{\mfa}{\faMobileAlt}

\usepackage{csquotes} % load it at the end to stop other packagings complaining